\section{Methods}

\subsection{Data}

We analyzed 12 cortical neurons from the MICrONS cortical connectome \citep{microns}: 6 excitatory (two layer 2/3 pyramidal, one layer 4 pyramidal, one layer 5 IT, one layer 5 ET, one layer 6 CT) and 6 inhibitory (two basket cells, two Martinotti cells, two bipolar cells). Dendritic morphologies were reconstructed as SWC files and synaptic inputs were extracted from the MICrONS HDF5 connectome database (v343), providing presynaptic neuron identity, 3D coordinates, and cell type classification for each synapse. Synapses were snapped to the nearest dendritic branch within 50~$\mu$m using the Neurostat framework \citep{neurostat}; synapses failing to snap were excluded ($< 1\%$ across all neurons). Valid synapse counts ranged from 488 to 2490 per neuron, with 8 to 107 presynaptic partners contributing $k \geq 5$ synapses each (Table~\ref{tab:summary}).

\subsection{Fast geodesic distance computation}

Pairwise geodesic distances along the dendritic tree are required for both the soma-distance null model and the per-partner clustering test. The existing implementation computes BFS per synapse pair, which is prohibitively slow for large distance matrices. We developed a fast geodesic distance computation that precomputes a matrix of shortest-path distances between all branch endpoints using Dijkstra's algorithm on the branch adjacency graph (each branch is an edge with weight equal to its total path length). Any pairwise synapse distance can then be computed in $O(1)$ by considering four possible paths through the two branch endpoints of each synapse's branch. For a neuron with $B$ branches and $N$ synapses, precomputation takes $O(B^2 \log B)$ and the full $N \times N$ distance matrix takes $O(N^2)$, reducing computation from ${\sim}30$~min to ${\sim}0.4$~s for $N = 1000$.

\subsection{Soma-distance null model}

To control for distance-dependent synapse density gradients, we constructed a null model in which synapse placement probability is proportional to the observed density profile rather than uniform along path length. The density profile was estimated as follows: (1)~compute the geodesic distance from the soma to each synapse; (2)~bin synapse counts by soma distance into 30 equal-width bins; (3)~compute the total dendritic path length available in each distance band by walking along all skeleton edges and accumulating edge lengths into the corresponding bin; (4)~compute raw density as synapse count divided by path length per bin; (5)~smooth with a Gaussian kernel (bandwidth = 10\% of the distance range). The dendritic tree was then divided into ${\sim}100$~nm segments, each assigned a weight proportional to segment length $\times$ density at its soma-distance midpoint. Mock synapse sets were generated by sampling segments proportional to these weights and placing synapses uniformly within each segment.

Spatial structure was tested using the multiscale variance-mean ratio (VMR) curves from the Neurostat framework \citep{neurostat}, compared against both the uniform null and the soma-distance null using Hotelling's $T^2$ test with 50 mock datasets per null model.

\subsection{Per-partner clustering test}

For each presynaptic partner contributing $k$ synapses to a postsynaptic neuron, we tested whether its synapses are more spatially compact than expected under a distance-matched permutation null. The test statistic was the mean pairwise geodesic distance among the partner's $k$ synapses, computed from the precomputed $N \times N$ distance matrix.

\subsubsection{Distance-matched permutation null}

The permutation null must preserve the partner's soma-distance distribution to avoid confounding distance gradients with clustering. For each permutation: (1)~all $N$ synapses on the neuron are binned into 10 equal-frequency soma-distance bins; (2)~the partner's $k$ synapses are profiled by their bin membership; (3)~$k$ synapses are resampled from all $N$ synapses, matching the partner's bin profile exactly; (4)~the mean pairwise geodesic distance is computed on the resampled set. This was repeated for 1000 permutations per partner. The $p$-value is the fraction of null values $\leq$ the observed value (one-sided test for clustering).

\subsubsection{Multiple testing correction}

Partners were tested at three synapse-count tiers ($k \geq 3$, $k \geq 5$, $k \geq 8$) to show how results change with statistical power. Within each tier, $p$-values were corrected using the Benjamini-Hochberg procedure at FDR $q = 0.05$. Bonferroni correction was also computed for comparison. Neuron-level enrichment was assessed by comparing the observed fraction of BH-significant partners to the expected fraction (0.05) using a one-sided binomial test.

\subsubsection{Compactness ratio}

To quantify the effect size for each partner, we defined the compactness ratio as:
\begin{equation}
    \text{compactness ratio} = \frac{\text{observed mean pairwise geodesic distance}}{\text{median of null mean pairwise geodesic distances}}
\end{equation}
Values below 1.0 indicate that the partner's synapses are more compact than expected under the distance-matched null.

\subsection{Label-shuffle control}

To verify that the clustering signal reflects partner-specific spatial organization rather than properties of the dendritic tree or the testing procedure, we performed a label-shuffle control. All synapse positions were held fixed, and presynaptic partner identity labels were randomly permuted across all synapses. The per-partner clustering test (at $k \geq 5$) was then re-run on the shuffled data. This was repeated for 1000 shuffles, with 200 permutations per partner per shuffle to manage computational cost. The fraction of BH-significant partners was recorded for each shuffle, and the real fraction was compared to the distribution of shuffled fractions.

\subsection{Software and reproducibility}

All analyses were performed in Python 3.9.6 using the Neurostat framework for dendritic spatial statistics and custom code for input-specific clustering analysis. Random seeds were fixed for all stochastic components (per-partner permutations: seed 42; null model generation: seeds 42--43). Code, execution instructions, and output checksums are provided in the accompanying reproducibility manifest.
